\chapter{Verificarea corectitudinii fragmentărilor realizate}

Pentru a demonstra că fragmentările realizate sunt corecte, trebuie să verificăm îndeplinirea a trei reguli fundamentale din teoria bazelor de date distribuite: regula completitudinii, regula reconstrucției și regula disjuncției.

\section*{1. Verificarea pentru fragmentarea orizontală (Tabela CLIENT)}

\begin{itemize}
    \item \textbf{Regula completitudinii:} Orice înregistrare din tabela globală \textit{CLIENT} se regăsește obligatoriu într-unul dintre cele două fragmente. Deoarece atributul de partiționare (\textit{tip\_client}) are o constrângere de tip CHECK care permite exclusiv valorile 'F' (Fizic) și 'J' (Juridic), orice client adăugat în sistem va fi repartizat automat în \textit{CLIENT\_FIZIC} sau \textit{CLIENT\_JURIDIC}. Nu există date omise sau "pierdute".
    \item \textbf{Regula reconstrucției:} Tabela globală inițială poate fi refăcută integral, fără alterări, reunind cele două fragmente prin operația de reuniune (UNION). Din punct de vedere matematic: 
    $$CLIENT = CLIENT\_FIZIC \cup CLIENT\_JURIDIC$$
    (La nivel de implementare practică, am asigurat această reconstrucție prin crearea view-urilor globale ce folosesc instrucțiunea \texttt{UNION ALL}).
    \item \textbf{Regula disjuncției:} Fragmentele obținute nu se suprapun (nu conțin înregistrări comune). Un client nu poate fi în același timp și persoană fizică, și persoană juridică, predicatele fiind reciproc exclusive. Astfel, intersecția fragmentelor este mulțimea vidă:
    $$CLIENT\_FIZIC \cap CLIENT\_JURIDIC = \emptyset$$
\end{itemize}

\section*{2. Verificarea pentru fragmentarea verticală (Tabela AGENT)}

\begin{itemize}
    \item \textbf{Regula completitudinii:} Toate coloanele (atributele) tabelei globale \textit{AGENT} au fost incluse în fragmentare. Atributele referitoare la profil se află în \textit{AGENT\_PROFIL}, iar atributele de securitate (email, parolă) se află în \textit{AGENT\_SEC}. Niciun atribut al relației inițiale nu a fost omis.
    \item \textbf{Regula reconstrucției:} Tabela globală originală poate fi recreată unind fragmentele pe baza cheii primare comune (\textit{agent\_id}), folosind operația de joncțiune naturală (JOIN). Matematic, acest lucru se exprimă astfel:
    $$AGENT = AGENT\_PROFIL \bowtie_{\text{agent\_id}} AGENT\_SEC$$
    (Această regulă a fost validată în implementare prin crearea view-ului \textit{v\_agent\_global}, care lipește cu succes profilul public de datele de login).
    \item \textbf{Regula disjuncției:} Cu excepția cheii primare (\textit{agent\_id}), care trebuie reprodusă obligatoriu în ambele fragmente pentru a face posibilă reconstrucția, nicio altă coloană nu este duplicată. Setul atributelor non-cheie din \textit{AGENT\_PROFIL} și setul atributelor non-cheie din \textit{AGENT\_SEC} sunt complet disjuncte.
\end{itemize}